\documentclass[12pt]{article}
\usepackage[utf8]{inputenc}
\usepackage[T1]{fontenc}
\usepackage[margin=1in]{geometry}
\usepackage{setspace}
\usepackage{xcolor}
\usepackage[colorlinks=true, urlcolor=blue!60!black, linkcolor=blue!60!black]{hyperref}
\setlength{\parindent}{2em}
\setlength{\parskip}{1.25em}
\renewcommand{\baselinestretch}{1.0}
\begin{document}
\begin{center}
{\LARGE \textit{Calder v. Bull}, 3 U.S. (3 Dall.) 386 (1798)\par}
\vspace{0.8em}
{\large Case Brief --- Four-Page Edition\par}
\end{center}
\vspace{1em}

\section{Memory Jogger}

A unanimous Court held that the Ex Post Facto Clause bars only retroactive criminal laws, not retroactive civil legislation, while Justices Chase and Iredell staged the founding American debate between natural-law constitutionalism and written-text positivism (\href{https://supreme.justia.com/cases/federal/us/3/386/}{Justia}).

\section{Detailed Case Facts}

In 1779 Normand Morrison executed a will leaving property to Caleb Bull and his wife. In March 1793 the Hartford County Court of Probate disapproved the will and refused to record it, so the estate passed by intestacy to Calder and his wife, Morrison's heiress. Connecticut law gave aggrieved parties 18 months to appeal; the Bulls let that period lapse, and no general statute authorized a new hearing.

Connecticut then had no modern written constitution separating powers; under its charter and long usage, the legislature historically exercised superintending judicial authority, including granting new trials. In May 1795 the legislature passed a resolution setting aside the 1793 decree and granting a new hearing. On rehearing the probate court approved the will, awarding the property to the Bulls. The Calders argued that the resolution was an \textit{ex post facto} law barred by Article I, \S~10, which provides that ``No State shall ... pass any ... ex post facto Law.''

\section{Procedural History}

The Superior Court at Hartford affirmed the new probate decree in 1796, and Connecticut's Supreme Court of Errors found no error, upholding the resolution under state law and the federal Constitution. The Calders brought a writ of error to the U.S. Supreme Court, which heard argument in February 1798 and decided the case on August 8, 1798.

\section{Judicial Votes}

Affirmed, \textbf{4--0}, by \textit{seriatim} opinions rather than a single opinion of the Court.
\begin{itemize}
\item \textbf{Seriatim opinions:} Justices Chase, Paterson, Iredell, and Cushing, all concurring in the judgment (Iredell expressly rejecting Chase's natural-law reasoning).
\item \textbf{Dissents:} None.
\item \textbf{Not participating:} Chief Justice Ellsworth and Justice Wilson.
\end{itemize}

\section{Holding}

Connecticut's 1795 resolution granting a new probate hearing did not violate the federal Constitution. The Ex Post Facto Clauses (Art. I, \S\S~9 and 10) reach only retroactive criminal laws, which Chase sorted into four categories: laws that criminalize conduct innocent when done, aggravate a crime, increase its punishment, or alter the rules of evidence to permit conviction on less testimony. Laws that ``mollify the rigor of the criminal law'' are not prohibited. The Court also has no jurisdiction to void a state statute for violating the state's own constitution. The Chase--Iredell clash over natural law was dicta, since both voted to affirm.

\section{Analysis of Opinions}

\textbf{Chase.} He could not ``subscribe to the omnipotence of a state legislature'': ``vital principles in our free republican governments'' forbid a legislature to make a man judge in his own cause or to take property from A and give it to B, because the people cannot be presumed to have entrusted such powers. On the Clause itself, however, drawing on English attainders, Blackstone, \textit{The Federalist}, and early state declarations of rights, he read ``ex post facto laws'' as a technical phrase limited to penal statutes. Reading it to cover civil laws would make the Contracts Clause and the Takings Clause ``unnecessary, and therefore improper.'' The 1793 decree, he added, had not vested a property right in Calder.

\textbf{Paterson.} Connecticut's legislature acted in a ``double capacity,'' so the resolution could be a customary judicial act. Even treated as legislation, ``ex post facto'' meant ``crimes, pains, and penalties.'' He confessed ``an ardent desire to have extended the provision in the Constitution to retrospective laws in general,'' but the text forbade it.

\textbf{Iredell.} Written constitutions define legislative limits; courts may void acts that violate them, but given the ``delicate and awful nature'' of that power, only in a ``clear and urgent case.'' Courts cannot void laws within constitutional power merely because judges think them contrary to natural justice, which has no fixed standard. Civil retroactivity, such as eminent domain for highways, is sometimes essential.

\textbf{Cushing.} The case was ``clear of all difficulty, taken either way'': judicial acts are untouched, and legislative ones are ``maintained and justified by the ancient and uniform practice of the State of Connecticut.''

\section{Examples: Future Applications}

\begin{enumerate}
\item \textbf{Same-Side (Permitted): Retroactive Estate-Tax Fix.} A retroactive amendment closing a tax loophole is a civil revenue measure, not punishment, though it remains subject to due-process review.
\item \textbf{Same-Side (Permitted): Felon License Bar.} A statute barring persons with prior felony convictions from holding a medical license is a nonpunitive civil regulation, not an \textit{ex post facto} law.
\item \textbf{Opposite-Side (Barred): Dropping Corroboration.} Eliminating a corroboration requirement for offenses already committed falls in Chase's fourth category (\textit{Carmell v. Texas}).
\item \textbf{Opposite-Side (Barred): Sentence Enhancement.} Retroactively increasing a mandatory prison term or extending a punitive restitution order falls in the third category (\textit{Ellingburg v. United States} (2026)).
\item \textbf{Fence-Sitter: Sex-Offender Registration.} A retroactive scheme labeled ``civil and regulatory'' escapes the Clause if genuinely nonpunitive (\textit{Smith v. Doe}), but banishment-like residency bars may cross into prohibited retroactive punishment.
\end{enumerate}

\section{Critique}

\textbf{Internal tension and progressive views.} Chase condemned laws taking property from A and giving it to B, yet upheld a resolution that shifted an inheritance from Calder to Bull on the technicality that no right had vested. Progressives admire Chase's refusal to accept legislative omnipotence as a precursor to unenumerated-rights protection, but warn that his vested-property examples foreshadowed \textit{Lochner}; most support the holding leaving legislatures leeway for retroactive civil legislation.

\textbf{Originalist views.} Originalists celebrate Iredell as the canonical statement of written constitutionalism, but some question the criminal-only holding: in \textit{Eastern Enterprises v. Apfel}, Justice Thomas doubted Chase's redundancy argument and signaled willingness to reconsider whether the Clauses reach retroactive civil laws.

\section{Key Quotations}

\begin{quote}
``1st. Every law that makes an action done before the passing of the law and which was innocent when done, criminal and punishes such action. 2d. Every law that aggravates a crime, or makes it greater than it was when committed. 3d. Every law that changes the punishment, and inflicts a greater punishment than the law annexed to the crime, when committed. 4th. Every law that alters the legal rules of evidence and receives less or different testimony than the law required at the time of the commission of the offence in order to convict the offender.''
\end{quote}
\noindent --- Justice Chase, 3 U.S. at 390.

\begin{quote}
``An act of the legislature (for I cannot call it a law) contrary to the great first principles of the social compact cannot be considered a rightful exercise of legislative authority.''
\end{quote}
\noindent --- Justice Chase, 3 U.S. at 388.

\begin{quote}
``Every ex post facto law must necessarily be retrospective, but every retrospective law is not an ex post facto law. The former only are prohibited.''
\end{quote}
\noindent --- Justice Chase, 3 U.S. at 391.

\end{document}
